\section{Class-number terms in the \(M_{11}\) multiplicities}
The class-number estimates apply to a specified contribution to a virtual
multiplicity; they do not bound a cusp contribution or force the
multiplicity to be positive.
We use the \(2\)-optimal virtual graded modules of
Cheng--Duncan--Mertens \cite[Theorem~4.2.5 and Sect.~4.3]{CDM2025},
with their nontrivial multiplier assignment at orders four and eight.
Let \(H_N(D)\) be their generalized Hurwitz class number for
\(\Gamma_0(N)\), and \(H=H_1\). Their normalization gives
\[
 H(D)=h(D)\quad(D<-4\text{ fundamental}),\qquad
 H(-3)=\tfrac13,\quad H(-4)=\tfrac12.
\]
This normalization must be fixed before applying the analytic
class-number formula.

For \(N=1,2,3,6\), let \(C_N(D)\) be the coefficient of \(2R_N\),
where
\[
 R_N=\frac{12}{\varphi(N)}
       \sum_{M\mid N}\mu(N/M)
       \frac{M}{[\mathrm{SL}_2(\Z):\Gamma_0(M)]}\,H_M.
\]
Here the equality denotes the corresponding generating functions,
or their coefficient at a fixed \(D\).
Define the class-number contribution by
\[
 A(D)=\frac{C_1(D)}{144}-\frac{C_2(D)}{48}
              +\frac{C_3(D)}{18}-\frac{C_6(D)}6.
\]

\begin{theorem}\label{thm:m11-class}
Let \(D=f^2D_0<0\) be a discriminant, with \(D_0\) fundamental,
\(f\ge1\), and \(3\nmid D\). Put \(q_0=|D_0|\) and
\begin{equation}\label{eq:m11-B}
 B(D)=\chi_{D_0}(3)\bigl(1-\chi_{D_0}(2)\bigr)
 \prod_{\substack{\ell^e\parallel f\\\ell\ {\nm odd}}}
 \left(1+(\ell-\chi_{D_0}(\ell))
                  \frac{\ell^e-1}{\ell-1}\right).
\end{equation}
Then
\[
 A(D)=\frac{\sqrt{q_0}}{\pi}\,B(D)L(1,\chi_{D_0}).
\]
For \(q\ge3\), define
\begin{align*}
 U(q)&=2e^\gamma
 \left(\log\log q-\log2+\frac12+\frac1{\log\log q}\right),\\
 V(q)&=\frac{\pi^2}{12e^\gamma}
 \left(\log\log q-\log2+\frac12+\frac1{\log\log q}
                  +\frac{14\log\log q}{\log q}\right)^{-1}.
\end{align*}
Alp\"oge's Dirichlet Riemann hypothesis and the explicit
Lamzouri--Li--Soundararajan bounds imply
\begin{equation}\label{eq:m11-bounds}
 \frac{|B(D)|\sqrt{q_0}}{\pi}V(q_0)
 \le |A(D)|\le
 \frac{|B(D)|\sqrt{q_0}}{\pi}U(q_0).
\end{equation}
If \(B(D)\ne0\), \(A(D)\) has its sign; if \(B(D)=0\), \(A(D)=0\).
For fundamental \(D\), this reduces to
\[
 A(D)=\chi_D(3)(1-\chi_D(2))H(D).
\]
\end{theorem}
\begin{proof}
Substitution of \(N=1,2,3,6\) gives
\[
 A=-2H_1+H_2+2H_3-H_6.
\]
Separate binary quadratic forms of discriminant \(D\) according to
their content \(d\mid f\), writing \(g=f/d\) and \(D'=g^2D_0\).
For a primitive form of discriminant \(D'\), the number of isotropic
lines modulo a prime \(\ell\) is \(1+\chi_{D'}(\ell)\).
If \(\ell\mid d\), the form after multiplication by \(d\) is zero
modulo \(\ell\), so all \(\ell+1\) lines are allowed.
Thus the contribution of this content to \(H_N(D)\), for squarefree
\(N\), is the weighted order class number \(H_{\nm prim}(D')\)
times
\[
 \prod_{\ell\mid N}\nu_\ell,\qquad
 \nu_\ell=
 \begin{cases}\ell+1&\ell\mid d,\\
 1+\chi_{D'}(\ell)&\ell\nmid d.
 \end{cases}
\]
This follows by counting the cosets of \(\Gamma_0(N)\), whose
local quotients are \(\mathbb P^1(\F_\ell)\); stabilizer weights
give the same formula for the two exceptional unit groups.
Consequently the multiplier in \(A\) is
\[
 -2+\nu_2+2\nu_3-\nu_2\nu_3=(\nu_3-1)(2-\nu_2).
\]
Since \(3\nmid D\), \(\nu_3-1=\chi_{D_0}(3)\).
The order class-number formula, including unit factors, is
\[
 H_{\nm prim}(g^2D_0)=
 \frac{\sqrt{q_0}}{\pi}L(1,\chi_{D_0})\,
 g\prod_{\ell\mid g}\left(1-\frac{\chi_{D_0}(\ell)}{\ell}\right).
\]
For every odd \(\ell^e\parallel f\), summing the local terms gives
\[
 1+\sum_{j=1}^e\ell^j(1-\chi_{D_0}(\ell)/\ell)
 =1+(\ell-\chi_{D_0}(\ell))\frac{\ell^e-1}{\ell-1}.
\]
For the prime two put \(a=\chi_{D_0}(2)\).
If \(2\nmid f\), the multiplier is \(1-a\).
If \(2^e\parallel f\), \(e\ge1\), the terms with
\(v_2(g)<e\) have multiplier \(-1\), and the last term has
multiplier \(1\). Their sum is
\[
 2^{e-1}(2-a)-
 \left(1+(2-a)\sum_{j=1}^{e-1}2^{j-1}\right)=1-a.
\]
This proves \eqref{eq:m11-B} and the exact formula for \(A\).
The analytic class-number formula is positive, so it also proves
the assertion about the sign. Finally
\cite[Theorem~1.5]{LLS2015} gives the explicit upper and lower
bounds for \(L(1,\chi)\) at sufficiently large conductor;
\cite[Theorem~1]{LT2022} extends these same bounds to all
primitive conductors \(q\ge3\). Applying these bounds proves
\eqref{eq:m11-bounds}.
\end{proof}

To specify the remaining terms for even discriminants, define
\[
 T_4(D)=\sum_{a^2+b^2+c^2=|D|/4}(-1)^{b+c},\qquad
 T_8(D)=\sum_{a^2+4b^2+4c^2=|D|/4}(-1)^{b+c},
\]
where the sums are over integers, and are zero if \(4\nmid D\).
Let \(a_D=C_{8|4}(D)\) be the coefficient of the normalized
cuspidal Jacobi form \(\varphi_{8|4}\) of
\cite[equation~(4.2.9)]{CDM2025}.
Absorb its fixed contribution in the base trace function into the
parameter \(\kappa\): if \(\lambda\) denotes the coefficient
parameter in \cite[equation~(4.3.1)]{CDM2025}, put
\(\kappa=\Re(\lambda)-4\).
The parameter lattice is unchanged by this translation.

\begin{proposition}\label{prop:m11-full}
With this choice of parameter the full multiplicity is
\[
 m_{55}(W_D)=A(D)+\frac14T_4(D)-\frac12T_8(D)
                       +\frac{\kappa}{4}a_D.
\]
For odd fundamental \(D\), this becomes
\[
 m_{55}(W_D)=
 \chi_D(3)(1-\chi_D(2))H(D)+\frac{\kappa}{4}a_D.
\]
\end{proposition}
\begin{proof}
The character inner product in
\cite[equation~(4.3.1)]{CDM2025} gives
\[
 m_{55}=A-\frac18C_4+\frac14C_8+
                          \frac{\Re(\lambda)}4a_D.
\]
The explicit trace functions in
\cite[equation~(4.2.11)]{CDM2025} give
\(C_4=-2T_4\) and \(C_8=-2T_8-4a_D\).
For example,
\(\theta^0_{4,0}-\theta^0_{4,4}
=\sum_{b\in\Z}(-1)^bq^{b^2}\), while the corresponding
difference for index \(16\) is
\(\sum_b(-1)^bq^{4b^2}\).
Multiplication by \(\theta^0_{1,0}\theta_{1,0}(\tau,z)\)
then gives exactly the two ternary sums above.
Substitution proves the first formula. For odd \(D\) the two
ternary sums vanish, and Theorem~\ref{thm:m11-class} gives the second.
\end{proof}

In particular, a two-sided bound for \(A(D)\) is a bound for
the multiplicity after these explicitly specified terms are subtracted.
When \(D\) is odd and fundamental, its class-number multiplier is
\(0\), \(2\), or \(-2\); reversing the sign reverses the order of
the bounds for the signed quantity.

\section{Central values and the size of a multiplicity}
Tunnell's formula determines the square of the coefficient in this
case, including the dependence on the twist; it does not determine
its sign or give a distribution law for the family.
Put
\[
 \frac{\eta(4\tau)^5\eta(16\tau)}
      {\eta(2\tau)^2\eta(8\tau)}
       =\sum_{n\ge1}a(n)q^n,\qquad
 \beta=\int_1^\infty\frac{dx}{\sqrt{x^3-x}}.
\]

\subsection{The exact normalization of the Tunnell coefficient}
\label{sec:tunnell-normalization}
Write
\[
 \frac{\eta(4\tau)^5\eta(16\tau)}{\eta(2\tau)^2\eta(8\tau)}
   =\sum_{n\geq1}a(n)q^n,
 \qquad
 \varphi_{8|4}(\tau,z)=\sum_{n,r}C_{\varphi}(r^2-4n)q^ny^r,
\]
where $\varphi_{8|4}=(y^{-1}+y)q+O(q^2)$ is the normalization in
Cheng--Duncan--Mertens \cite[equation~(4.2.9)]{CDM2025}.
The exact comparison in this normalization is
\begin{equation}\label{eq:tunnell-factor-two}
 C_{\varphi}(-n)=\frac12a(n)\qquad(n\equiv3\pmod8).
\end{equation}
In particular $C_{\varphi}(-3)=1$ whereas $a(3)=2$.
The factor two does not affect a nonvanishing criterion, but does affect
an identity of coefficients and any explicit constant multiplying their
squares.

\begin{proof}
First we identify the Jacobi form itself:
\begin{equation}\label{eq:phi84-explicit}
 \varphi_{8|4}(\tau,z)
   =\eta(2\tau)\eta(8\tau)^2\theta_{1,1}(\tau,z).
\end{equation}
Here is a modular verification, in addition to the leading-term check.
The Jacobi triple product gives
$\theta^0_{1,1}=2\eta(4\tau)^2/\eta(2\tau)$, so the right side is
\[
 2\eta(4\tau)^2\eta(8\tau)^2
    \frac{\theta_{1,1}(\tau,z)}{\theta^0_{1,1}(\tau)}.
\]
The eta criterion puts $\eta(4\tau)^2\eta(8\tau)^2$ in
$S_2(\Gamma_0(32))$ with trivial character: its two congruence sums are
$2\cdot4+2\cdot8=24$ and $2\cdot(32/4)+2\cdot(32/8)=24$,
and $4^2 8^2$ is a square.  The theta transformation law shows that
$\theta_{1,1}$ transforms by a scalar on $\Gamma_0(4)$, with the same
scalar at $z=0$; their ratio has weight zero and index one there.
It is holomorphic in the upper half-plane, since $\theta^0_{1,1}$ has
no zeros there.  Cusp holomorphy can be checked without dividing:
the positive eta product $\eta(2\tau)\eta(8\tau)^2$ has positive order
at every cusp, and the transformed index-one theta vector has only
nonnegative theta exponents.  Their product is therefore a Jacobi cusp
form of weight two and index one for $\Gamma_0(32)$.
Its first term is $(y^{-1}+y)q$.  The uniqueness specified in
Cheng--Duncan--Mertens \cite[equation~(4.2.9)]{CDM2025} proves \eqref{eq:phi84-explicit}.

Put $\vartheta(\tau)=\sum_{r\in\mathbb Z}q^{r^2}$.
Another triple-product identity gives
\[
 \frac{\eta(4\tau)^5\eta(16\tau)}{\eta(2\tau)^2\eta(8\tau)}
   =\vartheta(2\tau)\eta(8\tau)\eta(16\tau).
\]
The eta product on the right has exponents congruent to one modulo eight.
The terms with exponent congruent to three modulo eight come precisely
from the odd integers in $\vartheta(2\tau)$.  As
\[
 \sum_{r\ \mathrm{odd}}q^{2r^2}
    =2\frac{\eta(32\tau)^2}{\eta(16\tau)},
\]
we obtain the exact formal-series identity
\begin{equation}\label{eq:tunnell-sieve}
 \sum_{n\equiv3\ (8)}a(n)q^n=2\eta(8\tau)\eta(32\tau)^2.
\end{equation}
On the other hand the odd theta coefficient of
\eqref{eq:phi84-explicit}, with $\tau$ replaced by $4\tau$, is
\[
 \sum_{n\equiv3\ (8)} C_{\varphi}(-n)q^n
   =\eta(8\tau)\eta(32\tau)^2.
\]
Comparison with \eqref{eq:tunnell-sieve} proves
\eqref{eq:tunnell-factor-two} for every such $n$, without a finite
coefficient extrapolation.
\end{proof}

The coefficient $a(n)$ here is exactly the one in Tunnell's original
definition $g\vartheta_2$, where
$g=q\prod_{r\geq1}(1-q^{8r})(1-q^{16r})$
\cite[p.~323]{Tunnell1983}.  Thus the correction is a factor between the
two normalizations, not a change to Tunnell's definition.


\begin{theorem}\label{thm:tunnell-multiplicity}
Let \(D<0\) be squarefree with \(D\equiv21\pmod{24}\),
write \(n=|D|\), and let \(E_n:y^2=x^3-n^2x\).
For the parameter \(\kappa\) defined above one has
\[
 m_{55}(W_D)=\frac{\kappa}{8}a(n),\qquad
 L(E_n,1)=\frac{\beta}{4\sqrt n}\,a(n)^2.
\]
Consequently
\[
 m_{55}(W_D)^2=
       \frac{\kappa^2\sqrt n}{16\beta}L(E_n,1),\qquad
 |m_{55}(W_D)|=
       \frac{|\kappa|}{4\sqrt\beta}\,
       n^{1/4}\sqrt{L(E_n,1)}.
\]
If \(\kappa\ne0\), the zero multiplicities correspond exactly to
zero central values. For any \(X,M>0\), their exact counting
relation is
\[
 \begin{split}
 &\#\{n\le X:\ n\ {\nm squarefree},\ n\equiv3\ (24),\
                    |m_{55}(W_{-n})|\le M\}\\
 &\quad=
 \#\{n\le X:\ n\ {\nm squarefree},\ n\equiv3\ (24),\
             \sqrt n\,L(E_n,1)\le16\beta M^2/\kappa^2\}.
 \end{split}
\]
\end{theorem}
\begin{proof}
Here the discriminant is odd and divisible by three, so the
class-number multiplier in the direct calculation of
Proposition~\ref{prop:m11-full} is zero, and both ternary sums
vanish. The coefficient identity is \(a_D=a(n)/2\), as proved in
\eqref{eq:tunnell-factor-two}.
Tunnell's Theorem~3 \cite[pp.~328--329]{Tunnell1983}
gives the displayed central-value formula with this eta quotient
and this real period. Squaring the coefficient formula and
substituting gives the two identities. The counting relation is
their pointwise equivalence. If \(\kappa=0\), all these
multiplicities vanish and the claimed equivalence is inapplicable.
\end{proof}

A Riemann hypothesis for \(L(E_n,s)\) would place its nontrivial
zeros on \(\Re s=1\), which includes the central point itself.
It would also imply the conductor-aspect Lindel\"of estimate
\[
 L(E_n,1)\ll_\varepsilon n^\varepsilon,\qquad
 |m_{55}(W_{-n})|\ll_{\kappa,\varepsilon}n^{1/4+\varepsilon}.
\]
Indeed, the analytic conductor of the unitary normalization
$L(E_n,s+1/2)$ is $O(n^2)$, and the GRH upper bound of
Chandee \cite[Corollaries~1.1 and~1.3, p.~3]{Chandee2009}
is $C^{o(1)}$ as the conductor $C$ tends to infinity.
Substitution in the central-value identity gives the second estimate.
This is an upper bound, not a distribution theorem.
It would therefore neither force these multiplicities to be nonzero
nor determine their signs or their frequencies of occurrence.
The displayed counting identity is the precise transfer of any
independently proved central-value distribution to their absolute
values. A statement about ranks or congruent numbers requires its
own input beyond the location of zeros.
